% =====================================================================
% Chapter 5 — Results (Vysledky)
% Data Catalog thesis, Monika Pulcová
%
% Structure: timeline of the project across six semesters, followed by
% the final implemented features. No external citations here — all
% literature sits in the State of the Art and Discussion chapters.
%
% Paste into your Overleaf project under 05_vysledky.tex. Do not wrap
% in \chapter{} — your main.tex handles the chapter title.
% Figure files are referenced with placeholder names, adjust as needed.
% =====================================================================

\section{Timeline of the project}
\label{sec:results-timeline}

The project was developed across six semesters. Each semester had its own focus, and the application grew from a very small proof of concept into the full system described in this chapter. Old versions are still in the git history on separate branches.

\subsection{First semester: proof of concept in Plotly Dash}
\label{sec:results-timeline-s1}

The first semester was mostly about getting familiar with the data and with the basics of data visualization. We had a big text file with outputs of a machine learning model for lung cancer prediction. The file was hard to read for any user. So the goal was to parse it and show the data in a simple web application.

We used Python library Plotly Dash for this first version. Plotly Dash runs a minimal Flask server and allows to build a simple interactive frontend directly in Python with CSS. We rendered a heatmap of model predictions over the patients, with filtering by patient ID. There was no REST API, no user input, no database, just a local server that loaded the text file and drew the heatmap. The source code of this proof of concept is available in a separate repository.\footnote{\url{https://github.com/pulcovamon/semestral\_project\_I}}

This version was very simple, but it was the first step. It showed that the text file can be turned into something readable, and we got a first idea how to map procedure code sequences and predictions onto a 2D color grid.

\subsection{Second semester: first React and FastAPI application}
\label{sec:results-timeline-s2}

In the second semester the goal was to learn modern web development and to build a first proper web application with separate frontend and backend. A project of this size was completely new for us, so a lot of the semester was spent just learning React, FastAPI, REST API design, Docker, and basic deployment.

We implemented a simple FastAPI backend with a MySQL database for all the data. Celery with Redis was already used in this version for background tasks, because predictions had to run asynchronously. The API had a few basic endpoints to serve the catalog and to submit a sequence of codes for analysis. On top of that we built a React frontend with a minimal interface. Everything was packaged with Docker and Docker Compose, so the whole stack could be started with one command.

MySQL was used for everything, including the catalog data. The reason was simply that we had no experience with NoSQL databases at that point, so we stayed with what we knew. A lot of this code was later rewritten or thrown away, but the overall shape of the application (React frontend, FastAPI backend, Celery with Redis for background work, Docker Compose for deployment) is the same today. This version was presented to the supervisor and to other teachers at the end of the semester.

\subsection{Third semester: architecture, UX design, and model experiments}
\label{sec:results-timeline-s3}

The third semester was the planning semester. At this point we had a working but simple application, and it became clear that the architecture has to be properly designed before adding more features. So we made several documents.

The first one is the software components diagram (Fig.~\ref{fig:sw-components}). It shows the whole system split between Client and Server, with Nginx proxy as the entry point, REST API in the middle, and different storage engines for different purposes. In this design we moved away from the single MySQL database, because the data has different shapes and access patterns. The catalog data is hierarchical with lists of codes per patient, so a document database fits better. Medical codes with their descriptions are mostly queried by free-text search, so a search engine fits better. So in the end we have MySQL for users and model metadata, MongoDB for catalog data, Elasticsearch for medical codes and their definitions, and Redis for the Celery task queue and result backend.

The second document is the REST API design with endpoints for catalog access, prediction submission, results, and user management. The detailed port layout is shown in Fig.~\ref{fig:sw-ports} (Nginx on 80, frontend on 3000, catalog backend on 8081, prediction backend on 8080, MySQL on 3306 and 3307, Redis on 6379, web documentation on 8082).

The third document is the UX flowchart with three separate flows (Fig.~\ref{fig:ux-flowchart}): the registered user flow for doctors and RWE specialists, the anonymous scoring flow for quick predictions without login, and the catalog flow for browsing the training data. This design came from the course Návrh uživatelských rozhraní (User Interface Design). As part of that course we did user research for the two target groups (RWE specialist, clinical specialist), wrote user stories, made an HTA (hierarchical task analysis), and prepared a Figma mockup and prototype. The prototype was tested with four testers of different age, profession, and technical background. The test pointed to several issues, for example that users did not always follow the expected path for uploading a single patient or for interpreting the results. We used these findings for the final UX redesign.

In the same semester we also did the machine learning experiments. The idea was to train our own models with a simplified approach, just to have some real models to plug into the system. We started with HMM (Hidden Markov Model) as a natural choice for sequential data, but it did not work as we expected. The accuracy on our data was around 50\,\%, even on simulated data with clear patterns. So we concluded that HMM is probably not suitable for this task, at least not with the simplified setup, and moved to TF-IDF with classical classifiers. TF-IDF with logistic regression reached around 83\,\% accuracy, and TF-IDF with Random Forest reached up to 100\,\% on the held-out test set for the binary class. These experiments are available in the training folder of the main repository.\footnote{\url{https://github.com/pulcovamon/Disease-scoring-system/tree/main/training}} We then implemented parts of the new architecture in the application itself: split the databases, cleaned up the REST API, added login, and did the first UX redesign.

\subsection{Fourth to sixth semester: frontend redesign, AI-assisted development, polish}
\label{sec:results-timeline-s456}

The last three semesters were about finishing the system. We redesigned the frontend with help of AI agents, tried different agentic setups, and used them for UX tweaks, for translations between English and Czech, for small bug fixes, and for style consistency. AI-assisted frontend development was new for us, so the setup itself changed several times during these semesters.

The main changes in this period are the reworked data visualization (the catalog list, the accuracy heatmap, the patient detail view, and the color-coding modes described in Section~\ref{sec:results-catalog}), the multistep scoring form with CSV and JSON bulk upload, the start of ML model storage (not fully completed), light and dark mode, and the language switcher between English and Czech. We also switched from mock models to the actually uploaded models in the scoring backend, fixed many bugs that appeared in the process, added JSON input support next to CSV, and improved the medical code search.

At the moment some parts are still not working as they should. The model storage is only partially implemented, and several smaller bugs are still open. No new big features are planned, only bug fixing.

% -----------------------------------------------------------------
\begin{figure}[h]
    \centering
    % \includegraphics[width=\textwidth]{figures/sw_components.jpg}
    \caption{Software components diagram. The Client part contains the web application, the web documentation, and the Celery Flower monitoring. The Server part contains the REST API, three databases (MongoDB for catalog, MySQL for users and metadata), the Redis task queue, and the ML model worker. All HTTP traffic goes through the Nginx proxy.}
    \label{fig:sw-components}
\end{figure}
% -----------------------------------------------------------------

% -----------------------------------------------------------------
\begin{figure}[h]
    \centering
    % \includegraphics[width=0.9\textwidth]{figures/sw_ports.jpg}
    \caption{Port layout of the deployed system. External clients reach the Nginx proxy on port 80. Nginx routes to the React frontend (3000), the catalog backend (8081), the prediction backend (8080), and the generated web documentation (8082). The prediction backend talks to the Redis task queue (6379), the user database (MySQL, 3307), and to the Celery worker, which stores results back in Redis.}
    \label{fig:sw-ports}
\end{figure}
% -----------------------------------------------------------------

% -----------------------------------------------------------------
\begin{figure}[h]
    \centering
    % \includegraphics[width=\textwidth]{figures/ux_flowchart.jpg}
    \caption{UX flowchart with three user flows: the registered scoring flow for doctors and RWE specialists (left), the anonymous scoring flow for quick predictions without login (middle), and the catalog flow for browsing the training data (right).}
    \label{fig:ux-flowchart}
\end{figure}
% -----------------------------------------------------------------

\section{Training data catalog and longitudinal data visualization}
\label{sec:results-catalog}

Longitudinal patient data is hard to read in raw form. One patient can have hundreds of five-digit procedure codes, and the ordering, frequency, and clinical specialty of each code all matter. So the catalog offers three coordinated views of the same data. The first one is a list view, where each patient is a row with a short code preview and a model accuracy summary. The second one is an accuracy heatmap over the whole catalog, where each patient is a column and each scoring category is a row. The third one is a per-patient detail view with a smaller heatmap and the full code sequence with several color-coding options.

In the development database the catalog currently holds 22\,282 patient records. One patient has on average tens to hundreds of five-digit procedure codes, with big variation between patients. Estimated accuracy is defined here as the proportion of model predictions that agree with the ground-truth label stored together with the record. For example, patient 49 has 332 codes, 64 distinct codes, seven predictions made by different models, and estimated accuracy of 86\,\%. We use this patient as a reference case in the rest of the section, and most figures in this chapter show patient 49.

\subsection{Catalog list view}
\label{sec:results-catalog-list}

The list view (Fig.~\ref{fig:catalog-list}) is the default tab of the Training Data page. It shows one row per patient. Each row contains the patient ID, a preview of the first five procedure codes as rounded badges, a counter with the total code count (for example ``+41'', ``+480''), the number of predictions made for this patient, and the estimated accuracy averaged across all predictions for this patient. A ``Detail'' button on each row opens the patient detail view described in Section~\ref{sec:results-catalog-detail}.

The list is paginated. The default page size is 20 rows, and the user can change it. For the current database of 22\,282 records this gives 1\,115 pages, which is shown next to the page selector. A total-record counter is also kept in a pill element in the header of the page (``Total records: 22,282''), so the user knows immediately how big the dataset is.

Above the list there are two search inputs. The first one searches by Patient ID. The second one searches by a specific procedure code and highlights every cell where this code appears in the preview badges. So the user can scan which patients in the current cohort have a given code. The highlighting is done by synchronizing the search term with a URL query parameter, which is explained in Section~\ref{sec:results-catalog-url}.

% -----------------------------------------------------------------
\begin{figure}[h]
    \centering
    % \includegraphics[width=\textwidth]{figures/catalog_list.png}
    \caption{Catalog list view, first page of the training data catalog with total of 22\,282 records. Each row shows the patient ID, a preview of the first five procedure codes, the total code count, the number of predictions, and the estimated accuracy.}
    \label{fig:catalog-list}
\end{figure}
% -----------------------------------------------------------------

\subsection{Accuracy heatmap over the catalog}
\label{sec:results-catalog-heatmap}

The second tab on the Training Data page is the Accuracy Heatmap (Fig.~\ref{fig:catalog-heatmap}). It is a matrix where each column is one patient, identified by the patient number (1, 2, 3, \dots), and each row is one scoring category: ICD10 Multiclass, ICD10 Binary, and Active Phase. This view is made for the RWE analyst, who needs to scan many patients at once. The cell value is the per-patient accuracy of all predictions made in that scoring category, from 0\,\% to 100\,\%.

The cell color encodes the same number once again. We used a single-hue sequential purple colormap. 0\,\% accuracy gets a very light color close to white, and 100\,\% accuracy gets a deep purple. We used HSL color. Hue is 270 (purple), saturation 80\,\%, lightness goes from 85\,\% for 0\,\% accuracy down to 14.5\,\% for 100\,\%. So low and high values are easy to tell apart even in a dense view where many columns are next to each other. The text of the percentage is overlaid on the cell and automatically switches from black to white when the cell lightness drops below 50\,\%, so the number is always readable. The colormap is implemented without a charting library, only with the \texttt{d3-scale-chromatic} package and plain CSS. This keeps the JavaScript bundle small.

The heatmap uses the same pagination controls as the list view and shares the same URL state. So when a user filters the list to a specific patient ID or a specific code and then switches to the heatmap tab, the heatmap is narrowed down the same way. This way the two views stay synchronized.

% -----------------------------------------------------------------
\begin{figure}[h]
    \centering
    % \includegraphics[width=\textwidth]{figures/catalog_heatmap.png}
    \caption{Accuracy heatmap tab of the Training Data page. Columns are patients, rows are the three scoring categories (ICD10 Multiclass, ICD10 Binary, Active Phase). The single-hue purple palette encodes per-patient accuracy from 0\,\% (near white) to 100\,\% (deep purple).}
    \label{fig:catalog-heatmap}
\end{figure}
% -----------------------------------------------------------------

\subsection{Patient detail view}
\label{sec:results-catalog-detail}

When the user clicks ``Detail'' on a row in the list view, they are taken to the patient detail page (Fig.~\ref{fig:patient-detail}). At the top there is a row of four summary counters: total codes, distinct codes, number of predictions, and estimated accuracy. These four numbers show how long the history is, how many different codes it has, how many models ran on it, and their average accuracy.

Below the summary there is a smaller per-patient heatmap. It has the same three category rows as the catalog-level heatmap (ICD10 Multiclass, ICD10 Binary, Active Phase), but the columns now represent individual predictions made for this one patient, ordered in time (1, 2, 3, \dots). The color scale is different here. A binary green and red coding is used, where green means the model prediction agrees with the ground-truth label and red means it does not. So a red cell shows where the model was wrong for this patient, and the position of the red cell in the row says whether it happened early in the record or later.

% -----------------------------------------------------------------
\begin{figure}[h]
    \centering
    % \includegraphics[width=\textwidth]{figures/patient_detail.png}
    \caption{Patient detail view for patient 49 (332 codes, 64 distinct codes, 7 predictions, 86\,\% estimated accuracy). The upper part shows the per-patient agreement heatmap with binary green and red coding. The lower part shows the full procedure code sequence, here with the Specialty color-coding mode active.}
    \label{fig:patient-detail}
\end{figure}
% -----------------------------------------------------------------

\subsection{Color-coding modes for the code sequence}
\label{sec:results-catalog-colors}

Under the heatmap there is the full sequence of procedure codes for this patient, displayed as a grid of rounded code badges (Fig.~\ref{fig:patient-detail}, bottom part). For patient 49 this is 332 badges in chronological order. Raw badges alone are not easy to read, so four alternative color-coding modes are offered. The user switches between them with a tab row above the list (Specialty, TF-IDF (0), TF-IDF (1), Frequency). A legend below the tabs explains what the colors mean in the active mode.

The Specialty mode colors each badge by the medical specialty of the procedure (for example clinical oncology, pulmonology, physiology). Codes that cannot be mapped to a known specialty are shown in white with an ``Uncategorized'' legend entry. On hover, the full Czech name of the specialty is shown as a tooltip. So the user gets a sense of which department the patient interacted with at each visit.

The TF-IDF (0) and TF-IDF (1) modes visualize how informative each code is for a given class label, using the TF-IDF weight described in Section~\ref{sec:methods-tfidf}. Higher TF-IDF means the code is more specific for the selected label. Codes with a TF-IDF above a threshold are shown in a darker shade, so the clinician can visually pick out which procedures in the history most strongly drive the prediction toward label 0 or label 1. Both label 0 and label 1 have their own tab. So the user can switch between them and see which codes push the prediction one way and which push it the other way.

The Frequency mode colors each code by how many times this specific code appears across the whole training catalog. Rare codes are shown in light color, frequent codes in dark gray. So the user can see at a glance if the patient history is dominated by routine procedures or if it contains unusual ones.

Each of these modes uses a different sequential color scale from \texttt{d3-scale-chromatic}. The same color scale is used consistently inside one mode, so the user does not have to re-learn the legend for every patient. For the Specialty mode, which is categorical and not sequential, a qualitative palette is used instead.

\subsection{Filtering, pagination, and URL-based highlighting}
\label{sec:results-catalog-url}

All filter values in the catalog are stored in the browser URL, not only in component memory. For example, when the user searches for the code \texttt{89131}, the URL becomes \texttt{/catalog?code=89131} and every badge with value \texttt{89131} in the current page is highlighted. When the user then clicks ``Detail'' on patient 49, the \texttt{code} query parameter is passed along, so on the detail page the same code stays highlighted in the full 332-code sequence. And when the user clicks ``Back to Catalog'', the previous filter is restored. So the user does not lose their context.

This approach is useful for two reasons. A clinician or an RWE analyst can copy the current URL and send it to a colleague, and the colleague sees the exact same filtered view after login. So it is easy to share a specific cohort or a specific code of interest. Also the browser Back and Forward buttons work naturally inside the single-page application, because every filter change produces a new history entry.

The URL also stores pagination information. So a bookmark to page 42 of the heatmap view keeps the pagination offset and the current tab.

On top of the URL state there is client-side caching to avoid redundant backend calls. When the user moves from page 1 to page 2 and back to page 1, the first page is served from the cache. The same cache also serves the medical code dictionary used by the autocomplete in the scoring form. So the autocomplete responds instantly, without a backend query for every keystroke. The medical code dictionary is loaded once per session, which is enough because it rarely changes.